%% 3-16 データに直線を引く：最小二乗法（第12回：組 V3-8 が書いた）
%% 3-0 のばねのデータ（ソラさんの表）を回収する。数はすべて Python で計算して確かめた（v3-8w/d16.py）。
\makeatletter
\newcommand\bkthrqline[1]{\par\addvspace{3pt}{\centering #1\par}\addvspace{3pt}}
\newcommand\bkthrqstep[1]{\par\addvspace{6pt}\noindent{\sffamily\bfseries\color{bkDeep}#1}\par\nobreak\vspace{2pt}}
% 3-0 のソラさんのデータ（おもりの数、のび cm）
\def\bkthrqsora{(0,0) (1,1.6) (2,2.9) (3,4.6) (4,5.9) (5,7.6) (6,8.9)}
% ばねB（おもりの数、のび mm）
\def\bkthrqbane{(1,22) (2,39) (3,60)}
\makeatother

\bkchapter{3}{16}{データに\R{直線}{ちょく|せん}を\R{引}{ひ}く：\\ \R{最小二乗法}{さい|しょう|に|じょう|ほう}}{\bklevel{3}}{データに\R{直線}{ちょく|せん}を\R{引}{ひ}く：\R{最小二乗法}{さい|しょう|に|じょう|ほう}}
\bkmichishirube{3-0,3-5,3-8,3-15}{5-7,7-5}

\begin{goal}
\begin{itemize}[label={\color{\bkTT}●}, leftmargin=1.6\zw, itemsep=0pt, topsep=1pt]
\item \bklevel{3}\enspace ずれ（\R{残差}{ざん|さ}）の2\R{乗}{じょう}の\R{和}{わ}を\R{計算}{けい|さん}して、どの\R{直線}{ちょく|せん}がよく\R{合}{あ}うかをくらべられる。
\item \bklevel{3}\enspace \R{原点}{げん|てん}を\R{通}{とお}る\R{直線}{ちょく|せん} $y=ax$ の $a$ を、\R{平方完成}{へい|ほう|かん|せい}（3-8）で\R{求}{もと}められる。
\item \bklevel{3}\enspace \R{直線}{ちょく|せん} $y=ax+b$ の $a$ と $b$ を、\R{共分散}{きょう|ぶん|さん}と\R{分散}{ぶん|さん}から\R{求}{もと}め、\R{残差}{ざん|さ}で\R{確}{たし}かめられる。
\end{itemize}
\end{goal}

\begin{hatena}
\noindent\begin{minipage}[c]{0.52\linewidth}
3-0で、ソラさんがはかったばねの\R{点}{てん}に、3\R{本}{ぼん}の\R{直線}{ちょく|せん}を\R{引}{ひ}きました。$y=1.45x$、$y=1.5x$、$y=1.55x$ です。どの\R{線}{せん}も、\R{点}{てん}の\R{近}{ちか}くを\R{通}{とお}ります。\par
「\R{目}{め}で\R{見}{み}て、だいたい」では、\R{人}{ひと}によって\R{引}{ひ}く\R{線}{せん}が\R{変}{か}わってしまいます。だれが\R{計算}{けい|さん}しても\R{同}{おな}じ1\R{本}{ぽん}に\R{決}{き}まる、「いちばんよく\R{合}{あ}う\R{直線}{ちょく|せん}」の\R{決}{き}め\R{方}{かた}はあるでしょうか。3-0からの\R{宿題}{しゅく|だい}です。
\end{minipage}\hfill
\begin{minipage}[c]{0.45\linewidth}\centering
\begin{tikzpicture}
\begin{axis}[bkplot, width=52mm, height=48mm, xmin=0, xmax=6.8, ymin=0, ymax=10.5, xtick={0,...,6}, ytick={0,2,...,10}, xlabel={$x$（\R{個}{こ}）}, ylabel={$y$（cm）}]
  \addplot[bkOrange, line width=1.1pt, domain=0:6.5] {1.45*x};
  \addplot[bkBlue, line width=1.1pt, domain=0:6.5] {1.5*x};
  \addplot[bkGreen!80!black, line width=1.1pt, domain=0:6.5] {1.55*x};
  \addplot[only marks, mark=*, mark size=1.8pt, bkInk] coordinates {\bkthrqsora};
  \node[font=\sffamily\bfseries\Large, text=bkOrange!80!black] at (axis cs:1.6,8) {?};
\end{axis}
\end{tikzpicture}
\end{minipage}
\end{hatena}

3-0では、ばねにおもりを\R{下}{さ}げて、のびをはかりました。\R{表}{ひょう}にして、\R{点}{てん}をかいて、\R{定規}{じょう|ぎ}で\R{線}{せん}を\R{引}{ひ}くと、$y=1.5x$ という\R{式}{しき}ができました。でも、\R{点}{てん}は\R{線}{せん}の\R{上}{うえ}にぴったりのりません。はかった\R{数}{かず}には、かならずずれがあるからです（2-6）。

この\R{章}{しょう}では、3-0の\R{宿題}{しゅく|だい}に\R{答}{こた}えます。\R{使}{つか}う\R{道具}{どう|ぐ}は3つです。ずれを2\R{乗}{じょう}して\R{数}{かぞ}える\R{考}{かんが}え\R{方}{かた}（3-14の\R{分散}{ぶん|さん}と\R{同}{おな}じ）、2\R{次式}{じ|しき}の\R{平方完成}{へい|ほう|かん|せい}（3-8）、そして\R{共分散}{きょう|ぶん|さん}（3-15）です。これまで\R{学}{まな}んだことが、1つにつながります。

%% ---------------- やってみよう・思い出そう ----------------
\begin{yattemiyou}[\R{線}{せん}を\R{引}{ひ}いて、ずれをはかろう]
\noindent\begin{minipage}[c]{0.55\linewidth}
\R{別}{べつ}のばね（ばねB）に、おもりを1\R{個}{こ}、2\R{個}{こ}、3\R{個}{こ}と\R{下}{さ}げて、のびをmmではかりました。
\begin{bkcenter}
\begin{bktable}{colspec={Q[l,m,wd=26mm] *{3}{Q[c,m,wd=9mm]}}, rowsep=1pt, cells={font=\small}}
おもり $x$（\R{個}{こ}） & 1 & 2 & 3 \\
のび $y$（mm） & 22 & 39 & 60 \\
\end{bktable}
\end{bkcenter}
\begin{enumerate}[label={\bkcircled[bkGreen]{\arabic*}}, leftmargin=2.4\zw, itemsep=1pt, topsep=2pt]
\item \R{方眼}{ほう|がん}に3つの\R{点}{てん}をかこう。
\item 3\R{本}{ぼん}の\R{直線}{ちょく|せん} $y=18x$、$y=20x$、$y=22x$ を\R{引}{ひ}こう。どれも\R{原点}{げん|てん}を\R{通}{とお}ります。
\item それぞれの\R{線}{せん}について、\R{点}{てん}が\R{線}{せん}から\R{上下}{じょう|げ}にどれだけずれているかを、\R{表}{ひょう}にしよう。
\item どの\R{線}{せん}が、いちばんよく\R{合}{あ}っているかな？\quad「よく\R{合}{あ}う」を、\R{数}{かず}で\R{言}{い}えるかな？
\end{enumerate}
\end{minipage}\hfill
\begin{minipage}[c]{0.42\linewidth}\centering
\begin{tikzpicture}
\begin{axis}[bkplot, width=46mm, height=60mm, xmin=0, xmax=3.6, ymin=0, ymax=75, xtick={0,1,2,3}, ytick={0,10,...,70}, xlabel={$x$（\R{個}{こ}）}, ylabel={$y$（mm）}]
  \addplot[bkOrange, line width=1pt, domain=0:3.4] {18*x};
  \addplot[bkBlue, line width=1pt, domain=0:3.4] {20*x};
  \addplot[bkGreen!80!black, line width=1pt, domain=0:3.4] {22*x};
  \addplot[only marks, mark=*, mark size=2pt, bkInk] coordinates {\bkthrqbane};
  \node[font=\sffamily\footnotesize, text=bkGreen!50!black, anchor=south east, inner sep=1pt] at (axis cs:3.2,71) {$22x$};
  \node[font=\sffamily\footnotesize, text=bkBlue!80!black, anchor=west, inner sep=1pt] at (axis cs:3.42,66) {$20x$};
  \node[font=\sffamily\footnotesize, text=bkOrange!80!black, anchor=west, inner sep=1pt] at (axis cs:3.42,58) {$18x$};
\end{axis}
\end{tikzpicture}
\end{minipage}
\end{yattemiyou}

\begin{omoidasou}
\bkrecall{3-0}{ソラさんのばね：おもり0〜6\R{個}{こ}で、のびは 0、1.6、2.9、4.6、5.9、7.6、8.9\,cm。\R{目}{め}で\R{見}{み}て $y=1.5x$ を\R{引}{ひ}いた。}
\bkrecall{3-5}{\R{直線}{ちょく|せん} $y=ax+b$ の $a$ は\R{傾}{かたむ}き、$b$ は\R{切片}{せっ|ぺん}。\R{点}{てん} $(x_1,\ y_1)$ を\R{通}{とお}り\R{傾}{かたむ}き $m$ の\R{直線}{ちょく|せん}は $y-y_1=m(x-x_1)$。}
\bkrecall{3-8}{$(A-B)^2=A^2-2AB+B^2$。\R{平方完成}{へい|ほう|かん|せい}：$x^2+6x+10=(x+3)^2+1$。3-8の「もっと\R{知}{し}りたい」では、\R{点}{てん} $(1,\ 2.0)$、$(2,\ 4.3)$、$(3,\ 5.8)$ に $y=kx$ を\R{引}{ひ}き、$S=14(k-2)^2+0.13$ から $k=2$ を\R{決}{き}めた。}
\bkrecall{3-14}{\R{偏差}{へん|さ}＝\R{値}{あたい}－\R{平均}{へい|きん}。\R{偏差}{へん|さ}の2\R{乗}{じょう}の\R{平均}{へい|きん}が\R{分散}{ぶん|さん} $\sigma^2$。}
\bkrecall{3-15}{\R{共分散}{きょう|ぶん|さん} $\sigma_{xy}$ ＝（$x$ の\R{偏差}{へん|さ}×$y$ の\R{偏差}{へん|さ}）の\R{平均}{へい|きん}。}
\end{omoidasou}

\bkfukidashi{sora}{3-0では、\R{線}{せん}を\R{引}{ひ}くたびに\R{傾}{かたむ}きが\R{少}{すこ}しずつ\R{変}{か}わったんだ。\R{計算}{けい|さん}で1\R{本}{ぽん}に\R{決}{き}まるなら、すっきりするね。}

%% ---------------- 数学の道具 ----------------
\bksection{doguu}{\R{数学}{すう|がく}の\R{道具}{どう|ぐ}}

\bksubsection{1}{ずれを\R{数}{かず}で\R{表}{あらわ}す：\R{残差}{ざん|さ}と、その2\R{乗}{じょう}の\R{和}{わ}}

\leavevmode\label{bk:3-16:zure}やってみようの\R{線}{せん} $y=20x$ で\R{考}{かんが}えます。おもり1\R{個}{こ}のとき、\R{線}{せん}の\R{上}{うえ}の\R{値}{あたい}は $20\times1=20$ です。はかった\R{値}{あたい}は22なので、\R{点}{てん}は\R{線}{せん}より $22-20=2$ だけ\R{上}{うえ}にあります。このように、
\bkthrqline{ずれ＝（はかった\R{値}{あたい}）－（\R{線}{せん}の\R{値}{あたい}）}
とします。\R{点}{てん}が\R{線}{せん}より\R{上}{うえ}なら\R{正}{せい}、\R{下}{した}なら\R{負}{ふ}です。3つの\R{点}{てん}のずれは、\R{図}{ず}\ref{bk:3-16:fig-zure}のようになります。

\begin{bkfigure}
\begin{tikzpicture}
\begin{axis}[bkplot, width=100mm, height=68mm, xmin=0, xmax=3.5, ymin=0, ymax=70, xtick={0,1,2,3}, ytick={0,10,...,70}, xlabel={$x$（\R{個}{こ}）}, ylabel={$y$（mm）}]
  \addplot[bkBlue, line width=1.2pt, domain=0:3.4] {20*x};
  \addplot[bkRed, line width=1.6pt] coordinates {(1,20) (1,22)};
  \addplot[bkRed, line width=1.6pt] coordinates {(2,40) (2,39)};
  \addplot[only marks, mark=*, mark size=2.2pt, bkInk] coordinates {\bkthrqbane};
  \addplot[only marks, mark=o, mark size=2pt, bkBlue] coordinates {(1,20) (2,40) (3,60)};
  \node[font=\sffamily\footnotesize, text=bkRed!80!black, anchor=north west] at (axis cs:1.08,19) {ずれ $22-20=2$};
  \node[font=\sffamily\footnotesize, text=bkRed!80!black, anchor=west, align=left] at (axis cs:2.08,34) {ずれ $39-40=-1$};
  \node[font=\sffamily\footnotesize, text=bkRed!80!black, anchor=west, align=left] at (axis cs:3.08,55) {ずれ $60-60=0$};
  \node[font=\sffamily\footnotesize, text=bkBlue!80!black, anchor=south east] at (axis cs:3.3,67) {$y=20x$};
\end{axis}
\end{tikzpicture}
\bkfigcaption{\R{黒}{くろ}い\R{点}{てん}ははかった\R{値}{あたい}、\R{白}{しろ}い\R{丸}{まる}は\R{線}{せん}の\R{上}{うえ}の\R{値}{あたい}。\R{赤}{あか}い\R{線}{せん}\R{分}{ぶん}がずれ。}\label{bk:3-16:fig-zure}
\end{bkfigure}

\textbf{なぜ、たてのずれ？}\quad おもりの\R{数}{かず} $x$ は、\R{自分}{じ|ぶん}で\R{決}{き}めて\R{下}{さ}げるので、ずれはありません。ずれがあるのは、\R{目}{め}もりを\R{読}{よ}んだのび $y$ のほうです。だから、たての\R{方向}{ほう|こう}のずれを\R{考}{かんが}えます。\R{線}{せん}からはみ\R{出}{だ}して\R{残}{のこ}ったずれなので、これを\textbf{\R{残差}{ざん|さ}}ともいいます。\bkidx{ざんさ}{\R{残差}{ざん|さ}}

\bkthrqstep{ためし1　ずれをそのままたすと？}
$y=20x$ のずれをたすと、$2+(-1)+0=1$ です。\R{小}{ちい}さいので、よく\R{合}{あ}っていそうです。でも、ずれをそのままたすやり\R{方}{かた}には、\R{大}{おお}きな\R{落}{お}とし\R{穴}{あな}があります。\R{図}{ず}\ref{bk:3-16:fig-uchikeshi}を\R{見}{み}ましょう。

\begin{bkfigure}
\begin{tikzpicture}[x=1mm, y=1mm]
  \foreach \cx/\tt in {0/{よく\R{合}{あ}う\R{線}{せん}}, 68/{まったく\R{合}{あ}わない\R{線}{せん}}} {
    \draw[bkInk, line width=0.6pt, -{Stealth[length=3pt]}] (\cx,0) -- (\cx+56,0);
    \draw[bkInk, line width=0.6pt, -{Stealth[length=3pt]}] (\cx,0) -- (\cx,36);
    \node[font=\sffamily\small\bfseries, anchor=south] at (\cx+28,37.5) {\tt};}
  % よく合う線
  \draw[bkBlue, line width=1.2pt] (2,4) -- (54,32);
  \fill[bkInk] (16,11.5) circle (1.1) (40,24.5) circle (1.1);
  \node[font=\sffamily\footnotesize, anchor=north west] at (17,10.5) {ずれ0};
  \node[font=\sffamily\footnotesize, anchor=north west] at (41,23.5) {ずれ0};
  \node[font=\sffamily\footnotesize, text=bkBlue!80!black, anchor=north] at (28,-2) {たすと $0+0=0$};
  % 合わない線（右下がり）
  \draw[bkOrange, line width=1.2pt] (70,31) -- (122,5);
  \fill[bkInk] (84,11.5) circle (1.1) (108,24.5) circle (1.1);
  \draw[bkRed, line width=1.3pt] (84,11.5) -- (84,24);
  \draw[bkRed, line width=1.3pt] (108,24.5) -- (108,12);
  \node[font=\sffamily\footnotesize, text=bkRed!80!black, anchor=east] at (83,17.5) {$-12.5$};
  \node[font=\sffamily\footnotesize, text=bkRed!80!black, anchor=west] at (109,18) {$+12.5$};
  \node[font=\sffamily\footnotesize, text=bkOrange!80!black, anchor=north] at (96,-2) {たすと $(-12.5)+12.5=0$};
\end{tikzpicture}
\bkfigcaption{ずれをそのままたすと、まったく\R{合}{あ}わない\R{線}{せん}でも0になることがある。}\label{bk:3-16:fig-uchikeshi}
\end{bkfigure}

\R{右}{みぎ}の\R{線}{せん}は、\R{点}{てん}とは\R{逆}{ぎゃく}の\R{向}{む}きにかたむいていて、まったく\R{合}{あ}っていません。それなのに、\R{上}{うえ}へのずれ（$+12.5$）と\R{下}{した}へのずれ（$-12.5$）が\R{打}{う}ち\R{消}{け}し\R{合}{あ}って、\R{和}{わ}は0になってしまいます。これでは、よい\R{線}{せん}と\R{悪}{わる}い\R{線}{せん}を\R{見分}{み|わ}けられません。

\bkthrqstep{ためし2　ずれを2\R{乗}{じょう}してからたす}
そこで、3-14で\R{分散}{ぶん|さん}を\R{作}{つく}ったときと\R{同}{おな}じように、\textbf{ずれを2\R{乗}{じょう}してからたし}ます。2\R{乗}{じょう}すると、\R{負}{ふ}のずれも\R{正}{せい}になるので、\R{打}{う}ち\R{消}{け}し\R{合}{あ}いません。\R{図}{ず}\ref{bk:3-16:fig-uchikeshi}の\R{右}{みぎ}の\R{線}{せん}なら、$(-12.5)^2+12.5^2=312.5$ と\R{大}{おお}きな\R{数}{かず}になり、「\R{合}{あ}っていない」ことがわかります。さらに、2\R{乗}{じょう}すると\R{大}{おお}きなずれほどとても\R{大}{おお}きくなるので、1つの\R{点}{てん}から\R{大}{おお}きく\R{外}{そと}れた\R{線}{せん}は、きびしく\R{評価}{ひょう|か}されます。

この「ずれの2\R{乗}{じょう}の\R{和}{わ}」を、この\R{章}{しょう}では $S$ と\R{書}{か}くことにします。$S$ が\R{小}{ちい}さいほど、\R{線}{せん}は\R{点}{てん}によく\R{合}{あ}っています。やってみようの3\R{本}{ぼん}の\R{線}{せん}で\R{計算}{けい|さん}しましょう。

\begin{bkcenter}
\begin{bktable}{colspec={Q[c,m,wd=14mm] Q[c,m,wd=22mm] Q[c,m,wd=22mm] Q[c,m,wd=22mm] Q[c,m,wd=27mm]}, colsep=4pt, rowsep=1pt, cells={font=\small}, row{1}={font=\sffamily\bfseries\small}}
\R{線}{せん} & 1\R{個}{こ}のずれ & 2\R{個}{こ}のずれ & 3\R{個}{こ}のずれ & $S$（2\R{乗}{じょう}の\R{和}{わ}） \\
$y=18x$ & $22-18=4$ & $39-36=3$ & $60-54=6$ & $16+9+36=61$ \\
\SetCell{bg=bkYellow!30}$y=20x$ & \SetCell{bg=bkYellow!30}$2$ & \SetCell{bg=bkYellow!30}$-1$ & \SetCell{bg=bkYellow!30}$0$ & \SetCell{bg=bkYellow!30}$4+1+0=5$ \\
$y=22x$ & $22-22=0$ & $39-44=-5$ & $60-66=-6$ & $0+25+36=61$ \\
\end{bktable}
\bktabcaption{3\R{本}{ぼん}の\R{線}{せん}のずれと、ずれの2\R{乗}{じょう}の\R{和}{わ} $S$}
\end{bkcenter}

3\R{本}{ぼん}の\R{中}{なか}では、$S=5$ の $y=20x$ がいちばんよく\R{合}{あ}っています。では、$y=19.5x$ や $y=20.3x$ のような\R{線}{せん}なら、もっと\R{小}{ちい}さくなるかもしれません。$a$ をどんな\R{数}{かず}にしたら $S$ がいちばん\R{小}{ちい}さくなるかを、②で\R{計算}{けい|さん}で\R{決}{き}めます。

\begin{seikaku}
ずれの\R{大}{おお}きさをくらべるには、2\R{乗}{じょう}のかわりに、\R{符号}{ふ|ごう}をとった\R{大}{おお}きさ（\R{絶対値}{ぜっ|たい|ち}、3-2）をたす\R{方法}{ほう|ほう}もあります。でも、\R{絶対値}{ぜっ|たい|ち}の\R{和}{わ}をいちばん\R{小}{ちい}さくする\R{線}{せん}は、かんたんな\R{式}{しき}では\R{求}{もと}められません。2\R{乗}{じょう}の\R{和}{わ}なら、②・③のように\R{平方完成}{へい|ほう|かん|せい}で、ぴったり1つの\R{答}{こた}えが\R{式}{しき}で\R{出}{だ}せます。だから、ふつうは2\R{乗}{じょう}の\R{和}{わ}を\R{使}{つか}います。
\end{seikaku}

\bksubsection{2}{\R{原点}{げん|てん}を\R{通}{とお}る\R{直線}{ちょく|せん} $y=ax$：\R{平方完成}{へい|ほう|かん|せい}で $a$ を\R{決}{き}める}

\leavevmode\label{bk:3-16:genten}ばねは、おもりが0\R{個}{こ}ならのびも0です。だから、\R{引}{ひ}く\R{線}{せん}は\R{原点}{げん|てん}を\R{通}{とお}る $y=ax$ の\R{形}{かたち}にします（2-15の\R{比例}{ひ|れい}）。ばねBのデータで、$S$ がいちばん\R{小}{ちい}さくなる $a$ を、4つの\R{段}{だん}で\R{求}{もと}めましょう。

\bkthrqstep{\R{段}{だん}1　いろいろな $a$ で、$S$ を\R{計算}{けい|さん}してみる}
$a$ を18から22まで1ずつ\R{変}{か}えて、$S$ を\R{計算}{けい|さん}します（\R{計算}{けい|さん}のしかたは①の\R{表}{ひょう}と\R{同}{おな}じです）。

\noindent\begin{minipage}[c]{0.5\linewidth}\centering
\begin{bktable}{colspec={Q[c,m,wd=10mm] *{5}{Q[c,m,wd=7mm]}}, rowsep=1pt, colsep=3pt, cells={font=\small}}
$a$ & 18 & 19 & 20 & 21 & 22 \\
$S$ & 61 & 19 & \SetCell{bg=bkYellow!30}5 & 19 & 61 \\
\end{bktable}
\bktabcaption{$a$ を\R{変}{か}えたときの $S$}\par\vspace{2pt}
\raggedright \R{表}{ひょう}を\R{見}{み}ると、$a=20$ のまわりで $S$ が\R{左右}{さ|ゆう}\R{対称}{たい|しょう}に\R{変}{か}わっています。グラフにすると、\R{図}{ず}\ref{bk:3-16:fig-Sa}のように、\R{上}{うえ}に\R{開}{ひら}いた\R{放物線}{ほう|ぶつ|せん}（3-9、\R{谷}{たに}の\R{形}{かたち}）にならびます。いちばん\R{低}{ひく}い\R{点}{てん}は、$a=20$ のあたりにありそうです。
\end{minipage}\hfill
\begin{minipage}[c]{0.46\linewidth}\centering
\begin{tikzpicture}
\begin{axis}[bkplot, width=48mm, height=50mm, xmin=17, xmax=23, ymin=0, ymax=70, xtick={18,19,...,22}, ytick={0,20,...,60}, xlabel={$a$}, ylabel={$S$}, axis lines=left]
  \addplot[bkBlue, line width=1.2pt, domain=17.9:22.1, samples=60] {14*(x-20)^2+5};
  \addplot[only marks, mark=*, mark size=2pt, bkRed] coordinates {(18,61) (19,19) (20,5) (21,19) (22,61)};
\end{axis}
\end{tikzpicture}
\bkfigcaption{いちばん\R{低}{ひく}いのは $a=20$ で、$S=5$。}\label{bk:3-16:fig-Sa}
\end{minipage}

\bkthrqstep{\R{段}{だん}2　$a$ を\R{文字}{も|じ}のまま、$S$ を\R{式}{しき}にする}
1\R{個}{こ}めの\R{点}{てん}のずれは $22-a\times1$ です。その2\R{乗}{じょう}を、3-8の\R{式}{しき}
\[ (A-B)^2=A^2-2AB+B^2 \]
で\R{展開}{てん|かい}します。$A=22$、$B=a\times1$ と\R{考}{かんが}えると、
\[ (22-a\times1)^2=22^2-2\times a\times(1\times22)+a^2\times1^2 \]
どの\R{点}{てん}でも、$(y-ax)^2=y^2-2a\,xy+a^2x^2$ の\R{形}{かたち}です。$y^2$、$xy$、$x^2$ を\R{表}{ひょう}にして、たてにたします。

\begin{bkcenter}
\begin{bktable}{colspec={Q[c,m,wd=12mm] Q[c,m,wd=12mm] Q[c,m,wd=16mm] Q[c,m,wd=16mm] Q[c,m,wd=16mm]}, rowsep=0.5pt, cells={font=\small}, row{1}={font=\sffamily\bfseries\small}}
$x$ & $y$ & $x^2$ & $xy$ & $y^2$ \\
1 & 22 & 1 & 22 & 484 \\
2 & 39 & 4 & 78 & 1521 \\
3 & 60 & 9 & 180 & 3600 \\
\SetCell[c=2]{r} \R{合計}{ごう|けい} & & 14 & 280 & 5605 \\
\end{bktable}
\bktabcaption{$S$ を\R{式}{しき}にするための\R{表}{ひょう}}
\end{bkcenter}

3つの\R{点}{てん}の $(y-ax)^2$ をたすと、$a^2$ の\R{項}{こう}・$a$ の\R{項}{こう}・\R{数}{かず}の\R{項}{こう}に\R{分}{わ}けて\R{集}{あつ}められます。
\bkthrqline{$S=$（$y^2$ の\R{和}{わ}）$-\,2a\times$（$xy$ の\R{和}{わ}）$+\,a^2\times$（$x^2$ の\R{和}{わ}）}
\R{表}{ひょう}の\R{合計}{ごう|けい}を\R{入}{い}れると、
\begin{align*}
S&=5605-2a\times280+a^2\times14\\
&=14a^2-560a+5605
\end{align*}
$S$ は $a$ の2\R{次式}{じ|しき}です。\R{段}{だん}1の\R{表}{ひょう}で、\R{放物線}{ほう|ぶつ|せん}の\R{形}{かたち}にならんだのはこのためです。

\bkthrqstep{\R{段}{だん}3　\R{平方完成}{へい|ほう|かん|せい}する（3-8）}
$a^2$ の\R{係数}{けい|すう}14で、$a$ をふくむ2つの\R{項}{こう}をくくります（3-8の「\R{共通因数}{きょう|つう|いん|すう}でくくる」）。
\[ S=14(a^2-40a)+5605 \]
かっこの\R{中}{なか}を\R{平方完成}{へい|ほう|かん|せい}します。$a^2-40a$ は、$(a-20)^2=a^2-40a+400$ より400だけ\R{小}{ちい}さいので、
\[ a^2-40a=(a-20)^2-400 \]
これを\R{入}{い}れて、
\begin{align*}
S&=14\{(a-20)^2-400\}+5605\\
&=14(a-20)^2-5600+5605\\
&=14(a-20)^2+5
\end{align*}

\textbf{いちばん\R{小}{ちい}さくなるのはどこ？}\quad $(a-20)^2$ は2\R{乗}{じょう}なので、0\R{以上}{い|じょう}です。0になるのは $a=20$ のときだけです。だから $S$ は、$a=20$ のときにいちばん\R{小}{ちい}さく、そのときの\R{値}{あたい}は $14\times0+5=5$ です。\R{段}{だん}1の\R{表}{ひょう}やグラフとぴったり\R{合}{あ}います。$a=19.5$ なら $14\times0.25+5=8.5$、$a=20.3$ なら $14\times0.09+5=6.26$ で、どちらも5より\R{大}{おお}きくなります。

\bkthrqstep{\R{段}{だん}4　いつでも\R{使}{つか}える\R{形}{かたち}にする}
\R{段}{だん}2・\R{段}{だん}3を、どんなデータでもできるように\R{文字}{も|じ}で\R{書}{か}きます。$x^2$ の\R{和}{わ}を $P$、$xy$ の\R{和}{わ}を $Q$、$y^2$ の\R{和}{わ}を $R$ とすると、\R{段}{だん}2と\R{同}{おな}じように $S=Pa^2-2Qa+R$ です。\R{段}{だん}3と\R{同}{おな}じように\R{平方完成}{へい|ほう|かん|せい}すると、
\begin{align*}
S&=P\Bigl(a^2-\frac{2Q}{P}a\Bigr)+R\\
&=P\Bigl(a-\frac{Q}{P}\Bigr)^2+R-\frac{Q^2}{P}
\end{align*}
$S$ がいちばん\R{小}{ちい}さくなるのは、かっこの\R{中}{なか}が0になる $a=\dfrac{Q}{P}$ のときです。

\begin{taisetsu}
\R{原点}{げん|てん}を\R{通}{とお}る\R{直線}{ちょく|せん} $y=ax$ で、ずれの2\R{乗}{じょう}の\R{和}{わ}をいちばん\R{小}{ちい}さくする $a$ は\par
\hfill $a=$（$xy$ の\R{和}{わ}）÷（$x^2$ の\R{和}{わ}）\hfill\null
\end{taisetsu}

ばねBなら $a=280\div14=20$ で、\R{段}{だん}3と\R{同}{おな}じです。3-8の「もっと\R{知}{し}りたい」の3\R{点}{てん}なら、$x^2$ の\R{和}{わ}14、$xy$ の\R{和}{わ}28.0で、$28.0\div14=2$ です（3-8では\R{傾}{かたむ}きを $k$ と\R{書}{か}きました）。この\R{式}{しき}は\R{覚}{おぼ}えなくても、\R{段}{だん}2・\R{段}{だん}3の\R{平方完成}{へい|ほう|かん|せい}をすれば、いつでも\R{作}{つく}れます。

このように、ずれの2\R{乗}{じょう}の\R{和}{わ}をいちばん\R{小}{ちい}さくして\R{直線}{ちょく|せん}を\R{決}{き}める\R{方法}{ほう|ほう}を、\textbf{\R{最小二乗法}{さい|しょう|に|じょう|ほう}}といいます。\bkidx{さいしょうにじょうほう}{\R{最小二乗法}{さい|しょう|に|じょう|ほう}}「\R{二乗}{に|じょう}」は2\R{乗}{じょう}のことで、「2\R{乗}{じょう}の\R{和}{わ}を\R{最小}{さい|しょう}にする\R{方法}{ほう|ほう}」という\R{意味}{い|み}の\R{名前}{な|まえ}です。

\bkthrqstep{3-0の\R{宿題}{しゅく|だい}の\R{答}{こた}え}
\bkhatenaref の\R{答}{こた}えを\R{出}{だ}しましょう。ソラさんのデータで、$x^2$ と $xy$ を\R{表}{ひょう}にします（$x=0$ の\R{列}{れつ}は、$x^2$ も $xy$ も0なので\R{省}{はぶ}きます）。

\begin{bkcenter}
\begin{bktable}{colspec={Q[l,m,wd=16mm] *{6}{Q[c,m,wd=11mm]} Q[c,m,wd=14mm]}, rowsep=0.5pt, colsep=3pt, cells={font=\small}}
$x$ & 1 & 2 & 3 & 4 & 5 & 6 & \R{合計}{ごう|けい} \\
$y$ & 1.6 & 2.9 & 4.6 & 5.9 & 7.6 & 8.9 & \\
$x^2$ & 1 & 4 & 9 & 16 & 25 & 36 & 91 \\
$xy$ & 1.6 & 5.8 & 13.8 & 23.6 & 38.0 & 53.4 & 136.2 \\
\end{bktable}
\bktabcaption{3-0のソラさんのデータ}
\end{bkcenter}
\[ a=136.2\div91=1.4967\cdots \]
いちばんよく\R{合}{あ}う\R{直線}{ちょく|せん}は、$y=1.497x$、およそ $y=1.50x$ です。3-0でソラさんが\R{目}{め}で\R{見}{み}て\R{引}{ひ}いた $y=1.5x$ は、ほぼ\R{最小二乗法}{さい|しょう|に|じょう|ほう}の\R{答}{こた}えと\R{同}{おな}じでした。3\R{本}{ぼん}の\R{線}{せん}の $S$ をくらべると、$y=1.45x$ は0.258、$y=1.5x$ は0.060、$y=1.55x$ は0.318で、$y=1.5x$ がいちばん\R{小}{ちい}さくなります（\R{最小二乗法}{さい|しょう|に|じょう|ほう}の $y=1.497x$ では0.059）。だれが\R{計算}{けい|さん}しても、\R{同}{おな}じ1\R{本}{ぽん} $y=1.497x$ に\R{決}{き}まります。これで、\R{宿題}{しゅく|だい}の\R{答}{こた}えが\R{出}{で}ました。

\textbf{ばね\R{定数}{てい|すう}も\R{求}{もと}められる}\quad 3-0のおもりは1\R{個}{こ}\qty{20}{g}で、\R{重}{おも}さは\R{約}{やく}\qty{0.2}{N}です（\qty{100}{g}で\R{約}{やく}\qty{1}{N}、2-15）。1\R{個}{こ}あたり\qty{1.497}{cm}のびるので、
\[ \qty{1.497}{cm}\div\qty{100}{cm/m}=\qty{0.01497}{m},\qquad k=\frac{\qty{0.2}{N}}{\qty{0.01497}{m}}=\qty{13.4}{N/m} \]
ソラさんのばねのばね\R{定数}{てい|すう}は、\R{約}{やく}\qty{13}{N/m}です（2-15の $F=kx$）。

\bksubsection{3}{\R{原点}{げん|てん}を\R{通}{とお}らない\R{直線}{ちょく|せん} $y=ax+b$：\R{共分散}{きょう|ぶん|さん}÷\R{分散}{ぶん|さん}}

\leavevmode\label{bk:3-16:ippan}こんどは、のびではなく、ばねの\textbf{\R{全体}{ぜん|たい}の\R{長}{なが}さ}をはかった\R{場合}{ば|あい}です。\R{全体}{ぜん|たい}の\R{長}{なが}さは、おもりが0\R{個}{こ}でも、もとの\R{長}{なが}さがあるので、\R{比例}{ひ|れい}しません（2-15のまちがえやすいところ）。\R{線}{せん}は $y=ax+b$ の\R{形}{かたち}で、$a$ と $b$ の2つを\R{決}{き}めます。ばねCのデータです。

\begin{bkcenter}
\begin{bktable}{colspec={Q[l,m,wd=30mm] *{5}{Q[c,m,wd=10mm]}}, rowsep=1pt, cells={font=\small}}
おもり $x$（\R{個}{こ}） & 1 & 2 & 3 & 4 & 5 \\
\R{長}{なが}さ $y$（cm） & 14.0 & 15.9 & 18.1 & 20.1 & 21.9 \\
\end{bktable}
\bktabcaption{ばねCの\R{全体}{ぜん|たい}の\R{長}{なが}さ。\R{平均}{へい|きん}は $\bar{x}=3$、$\bar{y}=90.0\div5=18.0$}
\end{bkcenter}

\bkthrqstep{\R{段}{だん}1　\R{傾}{かたむ}きを\R{決}{き}めると、いちばんよい $b$ は「\R{平均}{へい|きん}」}
まず、\R{傾}{かたむ}きを $a=2$ と\R{決}{き}めたとして、いちばんよい $b$ を\R{探}{さが}します。ずれは $y-(2x+b)=(y-2x)-b$ です。$y-2x$ を\R{計算}{けい|さん}すると、
\[ 14.0-2=12.0,\qquad 15.9-4=11.9,\qquad 18.1-6=12.1, \]
\[ 20.1-8=12.1,\qquad 21.9-10=11.9 \]
です。この5つの\R{数}{かず}から $b$ をひいて2\R{乗}{じょう}した\R{和}{わ}を、いちばん\R{小}{ちい}さくすればよいのです。②と\R{同}{おな}じように\R{展開}{てん|かい}して\R{平方完成}{へい|ほう|かん|せい}すると（5つの\R{数}{かず}の\R{和}{わ}は60、2\R{乗}{じょう}の\R{和}{わ}は720.04）、
\begin{align*}
S&=5b^2-2\times60\times b+720.04\\
&=5(b^2-24b)+720.04\\
&=5\{(b-12)^2-144\}+720.04\\
&=5(b-12)^2-720+720.04=5(b-12)^2+0.04
\end{align*}
となり、$b=12$ のとき $S$ はいちばん\R{小}{ちい}さくなります。12は、5つの\R{数}{かず}12.0、11.9、12.1、12.1、11.9の\textbf{\R{平均}{へい|きん}}です。

これは、いつでも\R{成}{な}り\R{立}{た}ちます。いちばんよい $b$ は「$y-ax$ の\R{平均}{へい|きん}」、つまり $\bar{y}-a\bar{x}$ です。$b=\bar{y}-a\bar{x}$ を $y=ax+b$ に\R{入}{い}れると、$x=\bar{x}$ のとき $y=a\bar{x}+\bar{y}-a\bar{x}=\bar{y}$ です。つまり、\textbf{いちばんよい\R{直線}{ちょく|せん}は、\R{平均}{へい|きん}の\R{点}{てん} $(\bar{x},\ \bar{y})$ を\R{通}{とお}ります}（\R{図}{ず}\ref{bk:3-16:fig-heikin}）。

\begin{bkfigure}
\begin{tikzpicture}
\begin{axis}[bkplot, axis lines=left, width=96mm, height=62mm, xmin=0, xmax=6, ymin=10, ymax=26, xtick={0,...,5}, ytick={10,14,...,26}, xlabel={$x$（\R{個}{こ}）}, ylabel={$y$（cm）}]
  \addplot[bkGray, line width=0.8pt, dashed, domain=0.3:5.7] {18+1.2*(x-3)};
  \addplot[bkGray, line width=0.8pt, dashed, domain=0.3:5.7] {18+2.8*(x-3)};
  \addplot[bkBlue, line width=1.3pt, domain=0:5.7] {2*x+12};
  \addplot[only marks, mark=*, mark size=2pt, bkInk] coordinates {(1,14.0) (2,15.9) (3,18.1) (4,20.1) (5,21.9)};
  \addplot[only marks, mark=square*, mark size=2.4pt, bkRed] coordinates {(3,18)};
  \node[font=\sffamily\footnotesize, text=bkRed!80!black, anchor=north west, inner sep=2pt] at (axis cs:3.05,17.6) {\R{平均}{へい|きん}の\R{点}{てん} $(3,\ 18)$};
  \node[font=\sffamily\footnotesize, text=bkBlue!80!black, anchor=north west] at (axis cs:0.15,25.6) {\R{青}{あお}：$y=2x+12$};
  \node[font=\sffamily\footnotesize, text=bkGray, anchor=west] at (axis cs:1.3,11.2) {\R{点線}{てん|せん}：\R{傾}{かたむ}きのちがう\R{線}{せん}};
\end{axis}
\end{tikzpicture}
\bkfigcaption{よい\R{線}{せん}は、\R{平均}{へい|きん}の\R{点}{てん}を\R{中心}{ちゅう|しん}に\R{回}{まわ}して\R{探}{さが}せばよい。\R{点線}{てん|せん}も\R{平均}{へい|きん}の\R{点}{てん}を\R{通}{とお}るが、\R{傾}{かたむ}きが\R{合}{あ}っていない。}\label{bk:3-16:fig-heikin}
\end{bkfigure}

\bkthrqstep{\R{段}{だん}2　\R{平均}{へい|きん}の\R{点}{てん}を\R{通}{とお}る\R{線}{せん}で、\R{傾}{かたむ}きを\R{決}{き}める}
\R{平均}{へい|きん}の\R{点}{てん} $(3,\ 18)$ を\R{通}{とお}り、\R{傾}{かたむ}き $a$ の\R{直線}{ちょく|せん}は、3-5の\R{点}{てん}と\R{傾}{かたむ}きの\R{式}{しき}で
\[ y-18=a(x-3) \]
です。ずれは $(y-18)-a(x-3)$。これは「$y$ の\R{偏差}{へん|さ}」$-a\times$「$x$ の\R{偏差}{へん|さ}」です（3-14）。\R{偏差}{へん|さ}を\R{表}{ひょう}にします。

\begin{bkcenter}
\begin{bktable}{colspec={Q[c,m,wd=14mm] Q[c,m,wd=14mm] Q[c,m,wd=20mm] Q[c,m,wd=26mm] Q[c,m,wd=20mm]}, rowsep=0.5pt, cells={font=\small}, row{1}={font=\sffamily\bfseries\footnotesize}}
$x$ の\R{偏差}{へん|さ} & $y$ の\R{偏差}{へん|さ} & $x$ の\R{偏差}{へん|さ}の2\R{乗}{じょう} & \R{偏差}{へん|さ}の\R{積}{せき} & $y$ の\R{偏差}{へん|さ}の2\R{乗}{じょう} \\
$-2$ & $-4.0$ & 4 & 8.0 & 16.00 \\
$-1$ & $-2.1$ & 1 & 2.1 & 4.41 \\
0 & 0.1 & 0 & 0 & 0.01 \\
1 & 2.1 & 1 & 2.1 & 4.41 \\
2 & 3.9 & 4 & 7.8 & 15.21 \\
\SetCell[c=2]{r} \R{合計}{ごう|けい} & & 10 & 20.0 & 40.04 \\
\end{bktable}
\bktabcaption{ばねCの\R{偏差}{へん|さ}}
\end{bkcenter}
②の\R{段}{だん}2と\R{同}{おな}じように\R{展開}{てん|かい}すると（$y$ のかわりに $y$ の\R{偏差}{へん|さ}、$x$ のかわりに $x$ の\R{偏差}{へん|さ}）、
\[ S=10a^2-2\times20\times a+40.04=10(a-2)^2+0.04 \]
で、$a=2$ のときいちばん\R{小}{ちい}さくなります。

\bkthrqstep{\R{段}{だん}3　\R{共分散}{きょう|ぶん|さん}と\R{分散}{ぶん|さん}で\R{書}{か}く}
\R{段}{だん}2の $a$ は、（\R{偏差}{へん|さ}の\R{積}{せき}の\R{和}{わ}）÷（$x$ の\R{偏差}{へん|さ}の2\R{乗}{じょう}の\R{和}{わ}）$=20\div10=2$ です。\R{分子}{ぶん|し}と\R{分母}{ぶん|ぼ}を\R{両方}{りょう|ほう}とも\R{人数}{にん|ずう}5でわっても、\R{答}{こた}えは\R{変}{か}わりません。すると、\R{分子}{ぶん|し}は\R{偏差}{へん|さ}の\R{積}{せき}の\R{平均}{へい|きん}、つまり\textbf{\R{共分散}{きょう|ぶん|さん}} $\sigma_{xy}$（3-15）、\R{分母}{ぶん|ぼ}は $x$ の\R{偏差}{へん|さ}の2\R{乗}{じょう}の\R{平均}{へい|きん}、つまり $x$ の\textbf{\R{分散}{ぶん|さん}} $\sigma_x^{\,2}$（3-14）になります。

\begin{taisetsu}
$y=ax+b$ で、ずれの2\R{乗}{じょう}の\R{和}{わ}をいちばん\R{小}{ちい}さくする $a$ と $b$ は\par
\hfill $a=\dfrac{\sigma_{xy}}{\sigma_x^{\,2}}$（\R{共分散}{きょう|ぶん|さん}÷$x$ の\R{分散}{ぶん|さん}）\qquad $b=\bar{y}-a\bar{x}$\hfill\null\par
いちばんよい\R{直線}{ちょく|せん}は、\R{平均}{へい|きん}の\R{点}{てん} $(\bar{x},\ \bar{y})$ を\R{通}{とお}る。
\end{taisetsu}

ばねCでは、$\sigma_{xy}=20\div5=4$、$\sigma_x^{\,2}=10\div5=2$ で $a=4\div2=2$、$b=18-2\times3=12$。だから $y=2x+12$ です。$b=12$ は、おもりが0\R{個}{こ}のときの\R{長}{なが}さ、つまりばねのもとの\R{長}{なが}さが\R{約}{やく}\qty{12}{cm}だということです。はかっていない0\R{個}{こ}のときの\R{長}{なが}さまで、\R{直線}{ちょく|せん}からわかりました。

このように\R{最小二乗法}{さい|しょう|に|じょう|ほう}で\R{求}{もと}めた\R{直線}{ちょく|せん}を、\textbf{\R{回帰直線}{かい|き|ちょく|せん}}ともいいます。\bkidx{かいきちょくせん}{\R{回帰直線}{かい|き|ちょく|せん}}3-15のコラムのゴルトンがつけた\R{名前}{な|まえ}です。

\bkthrqstep{3-0の\R{表計算}{ひょう|けい|さん}の\R{式}{しき}のなぞ}
3-0では、\R{表計算}{ひょう|けい|さん}が $y=1.4893x+0.0321$ という\R{式}{しき}を\R{出}{だ}しました。これは、\R{原点}{げん|てん}を\R{通}{とお}ると\R{決}{き}めずに、$y=ax+b$ で\R{最小二乗法}{さい|しょう|に|じょう|ほう}をした\R{答}{こた}えです。ソラさんのデータは、$\bar{x}=3$、$\bar{y}=31.5\div7=4.5$、\R{偏差}{へん|さ}の\R{積}{せき}の\R{和}{わ}が41.7、$x$ の\R{偏差}{へん|さ}の2\R{乗}{じょう}の\R{和}{わ}が28なので、
\[ a=41.7\div28=1.4893\cdots,\qquad b=4.5-1.4893\times3=0.0321\cdots \]
ぴったり\R{合}{あ}いました。\R{表計算}{ひょう|けい|さん}の\R{中}{なか}でも、この\R{章}{しょう}と\R{同}{おな}じ\R{計算}{けい|さん}をしていたのです。\R{表計算}{ひょう|けい|さん}では、\texttt{=SLOPE(B2:B8,A2:A8)} で\R{傾}{かたむ}き、\texttt{=INTERCEPT(B2:B8,A2:A8)} で\R{切片}{せっ|ぺん}を\R{求}{もと}められます（はじめに $y$ の\R{範囲}{はん|い}、\R{次}{つぎ}に $x$ の\R{範囲}{はん|い}）。

\textbf{どちらの\R{式}{しき}を\R{使}{つか}う？}\quad ②の $y=1.497x$ と、この $y=1.4893x+0.0321$ は、\R{少}{すこ}しちがいます。おもりが0\R{個}{こ}ならのびは0、という\R{物理}{ぶつ|り}のきまりを\R{知}{し}っているなら、\R{原点}{げん|てん}を\R{通}{とお}る $y=ax$ を\R{使}{つか}います。もとの\R{長}{なが}さのように、0\R{個}{こ}のときの\R{値}{あたい}がわからないなら、$y=ax+b$ を\R{使}{つか}います。どちらの\R{形}{かたち}の\R{式}{しき}を\R{使}{つか}うかは、\R{数学}{すう|がく}ではなく、しくみの\R{理解}{り|かい}で\R{決}{き}めるのです。

\bksubsection{4}{\R{残差}{ざん|さ}を\R{見}{み}て、\R{直線}{ちょく|せん}でよいか\R{確}{たし}かめる}

\leavevmode\label{bk:3-16:zansa}\R{最小二乗法}{さい|しょう|に|じょう|ほう}は、どんなデータにも、いちばん\R{合}{あ}う\R{直線}{ちょく|せん}を1\R{本}{ぽん}\R{出}{だ}してくれます。\R{点}{てん}が\R{曲}{ま}がってならんでいても、\R{直線}{ちょく|せん}を\R{出}{だ}してしまいます。だから\R{最後}{さい|ご}に、\textbf{\R{残差}{ざん|さ}（ずれ）のならび\R{方}{かた}}を\R{見}{み}て、\R{直線}{ちょく|せん}でよいかを\R{確}{たし}かめます。

\noindent\begin{minipage}[c]{0.5\linewidth}\centering
\begin{tikzpicture}
\begin{axis}[bkplot, axis lines=left, width=62mm, height=46mm, xmin=0, xmax=6.5, ymin=-0.2, ymax=0.2, xtick={0,...,6}, ytick={-0.2,-0.1,0,0.1,0.2}, xlabel={$x$}, ylabel={\R{残差}{ざん|さ}（cm）}, yticklabel style={/pgf/number format/fixed}, title={\sffamily\small ばね（3-0）}]
  \addplot[bkInk, line width=0.7pt, domain=0:6.5] {0};
  \addplot[ycomb, bkBlue, line width=1.4pt, mark=*, mark size=1.6pt] coordinates {(0,0) (1,0.103) (2,-0.093) (3,0.110) (4,-0.087) (5,0.117) (6,-0.080)};
\end{axis}
\end{tikzpicture}
\bkfigcaption{0のまわりに、＋と－がまざる。}\label{bk:3-16:fig-zansa1}
\end{minipage}\hfill
\begin{minipage}[c]{0.48\linewidth}\centering
\begin{tikzpicture}
\begin{axis}[bkplot, axis lines=left, width=62mm, height=46mm, xmin=0, xmax=28, ymin=-0.8, ymax=0.8, xtick={0,5,...,25}, ytick={-0.5,0,0.5}, xlabel={$x$}, ylabel={\R{残差}{ざん|さ}（cm）}, title={\sffamily\small \R{輪}{わ}ゴム}]
  \addplot[bkInk, line width=0.7pt, domain=0:28] {0};
  \addplot[ycomb, bkRed, line width=1.4pt, mark=*, mark size=1.6pt] coordinates {(5,0.5) (10,-0.25) (15,-0.5) (20,-0.25) (25,0.5)};
\end{axis}
\end{tikzpicture}
\bkfigcaption{＋、－、－、－、＋の\R{谷}{たに}の\R{形}{かたち}。}\label{bk:3-16:fig-zansa2}
\end{minipage}

\R{図}{ず}\ref{bk:3-16:fig-zansa1}は、3-0のソラさんのデータと $y=1.497x$ の\R{残差}{ざん|さ}です。どれも0.12\,cmより\R{小}{ちい}さく、＋と－がまざっていて、きまったならび\R{方}{かた}がありません。これは、はかるときの\R{小}{ちい}さなずれ（2-6）だと\R{考}{かんが}えられます。\R{直線}{ちょく|せん}でよいということです。

\R{図}{ず}\ref{bk:3-16:fig-zansa2}は、3-0のやってみようの\R{輪}{わ}ゴムで、10\R{円玉}{えん|だま}を5まい、10まい、…と\R{入}{い}れたときののび（1.0、2.0、3.5、5.5、8.0\,cm）に、\R{最小二乗法}{さい|しょう|に|じょう|ほう}で\R{直線}{ちょく|せん} $y=0.35x-1.25$ を\R{引}{ひ}いたときの\R{残差}{ざん|さ}です。\R{両}{りょう}はしが＋、まん\R{中}{なか}が－の、\R{谷}{たに}の\R{形}{かたち}にならんでいます。\R{点}{てん}が\R{直線}{ちょく|せん}ではなく、\R{曲}{ま}がったきまりにしたがっているしるしです。\R{輪}{わ}ゴムは、のびるにつれてのび\R{方}{かた}が\R{変}{か}わるのでした（3-0）。このようなときは、\R{直線}{ちょく|せん}ではなく、\R{曲}{ま}がった\R{式}{しき}を\R{考}{かんが}えます（\R{第}{だい}2\R{部}{ぶ}の3-7〜3-10）。

\begin{taisetsu}
\R{残差}{ざん|さ}が、0のまわりに\R{小}{ちい}さく、＋と－がばらばらにならぶ → \R{直線}{ちょく|せん}でよい\par
\R{残差}{ざん|さ}が、\R{山}{やま}や\R{谷}{たに}の\R{形}{かたち}にならぶ → \R{曲}{ま}がった\R{関係}{かん|けい}。\R{直線}{ちょく|せん}は\R{合}{あ}わない\par
\R{残差}{ざん|さ}が1つだけとても\R{大}{おお}きい → はかりまちがいでないか、\R{確}{たし}かめる
\end{taisetsu}

\begin{motto}[3]
\textbf{\R{微分}{び|ぶん}で\R{最小}{さい|しょう}を\R{求}{もと}める（5-7）}\par
②では、$S$ がいちばん\R{小}{ちい}さくなる $a$ を、\R{平方完成}{へい|ほう|かん|せい}で\R{求}{もと}めました。\R{第}{だい}5\R{巻}{かん}の5-7では、「グラフのいちばん\R{低}{ひく}い\R{点}{てん}では、グラフの\R{傾}{かたむ}きが0になる」ことを\R{使}{つか}って、いちばん\R{小}{ちい}さい\R{所}{ところ}を\R{求}{もと}める\R{方法}{ほう|ほう}（\R{微分}{び|ぶん}）を\R{学}{まな}びます。$a$ と $b$ の2つを\R{同}{おな}\R{時}{じ}に\R{決}{き}めるときも、たくさんの\R{式}{しき}を\R{同}{おな}\R{時}{じ}に\R{解}{と}けば\R{求}{もと}められます（\R{行列}{ぎょう|れつ}、7-5）。\par\vspace{4pt}
\textbf{\R{曲}{ま}がったデータも、\R{直線}{ちょく|せん}にして\R{最小二乗法}{さい|しょう|に|じょう|ほう}（3-12）}\par
\R{放射性物質}{ほう|しゃ|せい|ぶっ|しつ}の\R{量}{りょう}や、\R{冷}{さ}めていくお\R{湯}{ゆ}の\R{温度}{おん|ど}\R{差}{さ}のように、\R{倍々}{ばい|ばい}でへるデータ（3-11）は、たての\R{軸}{じく}の\R{目}{め}もりを\R{片対数}{かた|たい|すう}グラフ（3-12）にすると、\R{点}{てん}が\R{直線}{ちょく|せん}にならびます。その\R{直線}{ちょく|せん}に\R{最小二乗法}{さい|しょう|に|じょう|ほう}を\R{使}{つか}えば、\R{半分}{はん|ぶん}になる\R{時間}{じ|かん}などを、データから\R{決}{き}められます。
\end{motto}

\begin{kotoba}
\small
\bkword{\R{残差}{ざん|さ}}{はかった\R{値}{あたい}－\R{直線}{ちょく|せん}の\R{値}{あたい}。\R{点}{てん}が\R{直線}{ちょく|せん}から\R{上下}{じょう|げ}にずれた\R{大}{おお}きさ。}
\bkword{\R{最小二乗法}{さい|しょう|に|じょう|ほう}}{ずれ（\R{残差}{ざん|さ}）の2\R{乗}{じょう}の\R{和}{わ}がいちばん\R{小}{ちい}さくなるように、\R{直線}{ちょく|せん}を\R{決}{き}める\R{方法}{ほう|ほう}。}
\bkword{\R{回帰直線}{かい|き|ちょく|せん}}{\R{最小二乗法}{さい|しょう|に|じょう|ほう}で\R{求}{もと}めた\R{直線}{ちょく|せん}。\R{平均}{へい|きん}の\R{点}{てん} $(\bar{x},\ \bar{y})$ を\R{通}{とお}る。}
\end{kotoba}

%% ---------------- 例題 ----------------
\bksection{reidai}{\R{例題}{れい|だい}でたしかめよう}

\begin{reidai}[ずれの2\R{乗}{じょう}の\R{和}{わ}で、2\R{本}{ほん}の\R{線}{せん}をくらべる]
3つの\R{点}{てん} $(1,\ 3)$、$(2,\ 4)$、$(3,\ 7)$ があります。\R{直線}{ちょく|せん} $y=2x$ と $y=2.5x$ のどちらがよく\R{合}{あ}っているかを、ずれの2\R{乗}{じょう}の\R{和}{わ} $S$ でくらべましょう。\par
\kangaekata それぞれの $x$ で「はかった\R{値}{あたい}－\R{線}{せん}の\R{値}{あたい}」を\R{計算}{けい|さん}し、2\R{乗}{じょう}してたします。$S$ が\R{小}{ちい}さいほうが、よく\R{合}{あ}っています。\par
\shiki $y=2x$：ずれは $3-2=1$、$4-4=0$、$7-6=1$\par
\hspace*{2.6\zw}$S=1^2+0^2+1^2=2$\par
\hspace*{2.6\zw}$y=2.5x$：ずれは $3-2.5=0.5$、$4-5=-1$、$7-7.5=-0.5$\par
\hspace*{2.6\zw}$S=0.25+1+0.25=1.5$\par
\kotae $y=2.5x$ のほうがよく\R{合}{あ}っている（$S$ が2と1.5で、1.5のほうが\R{小}{ちい}さい）
\end{reidai}

\begin{column}{\R{最小二乗法}{さい|しょう|に|じょう|ほう}で、\R{星}{ほし}を\R{見}{み}つけた}
1801\R{年}{ねん}のお\R{正月}{しょう|がつ}、イタリアの\R{天文学者}{てん|もん|がく|しゃ}が、\R{火星}{か|せい}と\R{木星}{もく|せい}のあいだに\R{新}{あたら}しい\R{星}{ほし}（\R{小惑星}{しょう|わく|せい}ケレス）を\R{見}{み}つけました。ところが\R{星}{ほし}はすぐに\R{太陽}{たい|よう}のまぶしさにかくれ、\R{見}{み}えなくなりました。はかった\R{位置}{い|ち}には\R{少}{すこ}しずつずれがあり、どこにもどってくるのか、だれにもわかりませんでした。\par
ドイツの\R{若}{わか}い\R{数学者}{すう|がく|しゃ}ガウスは、ずれのあるデータから「ずれの2\R{乗}{じょう}の\R{和}{わ}がいちばん\R{小}{ちい}さくなる\R{道}{みち}すじ」を\R{計算}{けい|さん}して、\R{星}{ほし}がもどってくる\R{場所}{ば|しょ}を\R{予想}{よ|そう}しました。その\R{年}{ねん}の\R{終}{お}わり、\R{望遠鏡}{ぼう|えん|きょう}を\R{向}{む}けると、ケレスはほぼ\R{予想}{よ|そう}どおりの\R{場所}{ば|しょ}に\R{見}{み}つかりました。こうして\R{最小二乗法}{さい|しょう|に|じょう|ほう}は、\R{世界}{せ|かい}に\R{知}{し}られるようになりました。
\end{column}

\begin{reidai}[\R{原点}{げん|てん}を\R{通}{とお}る\R{直線}{ちょく|せん}の $a$ を、\R{平方完成}{へい|ほう|かん|せい}で\R{求}{もと}める]
3つの\R{点}{てん} $(1,\ 4)$、$(2,\ 7)$、$(3,\ 8)$ に\R{直線}{ちょく|せん} $y=ax$ を\R{引}{ひ}きます。ずれの2\R{乗}{じょう}の\R{和}{わ} $S$ を $a$ の\R{式}{しき}にして\R{平方完成}{へい|ほう|かん|せい}し、$S$ がいちばん\R{小}{ちい}さくなる $a$ を\R{求}{もと}めましょう。\par
\kangaekata $x^2$、$xy$、$y^2$ の\R{和}{わ}を\R{求}{もと}めて、②の\R{段}{だん}2と\R{同}{おな}じように
\bkthrqline{$S=$（$x^2$ の\R{和}{わ}）$a^2-2$（$xy$ の\R{和}{わ}）$a+$（$y^2$ の\R{和}{わ}）}
を\R{作}{つく}ります。\par
\shiki $x^2$ の\R{和}{わ}は14、$xy$ の\R{和}{わ}は $4+14+24=42$、$y^2$ の\R{和}{わ}は $16+49+64=129$\par
\hspace*{2.6\zw}$S=14a^2-84a+129$\par
\hspace*{2.6\zw}$\phantom{S}=14(a^2-6a)+129$\quad 14でくくる\par
\hspace*{2.6\zw}$\phantom{S}=14\{(a-3)^2-9\}+129$\quad $a^2-6a=(a-3)^2-9$\par
\hspace*{2.6\zw}$\phantom{S}=14(a-3)^2+3$\par
\kotae $a=3$（そのとき $S=3$ でいちばん\R{小}{ちい}さい）。\R{確}{たし}かめ：$42\div14=3$
\end{reidai}

\begingroup\tcbset{equal height group=bkthrqjA}%
\noindent\begin{minipage}[t]{0.49\linewidth}
\begin{jibun}
3つの\R{点}{てん} $(1,\ 2)$、$(2,\ 5)$、$(3,\ 6)$ があります。\R{直線}{ちょく|せん} $y=2x$ と $y=2.5x$ のどちらがよく\R{合}{あ}っているかを、ずれの2\R{乗}{じょう}の\R{和}{わ}でくらべましょう。
\bkjans{$y=2x$：ずれ0、1、0で $S=1$。$y=2.5x$：ずれ $-0.5$、0、$-1.5$ で $S=2.5$。$y=2x$ のほうがよく\R{合}{あ}う}
\end{jibun}
\end{minipage}\hfill
\begin{minipage}[t]{0.49\linewidth}
\begin{jibun}
3つの\R{点}{てん} $(1,\ 3)$、$(2,\ 5)$、$(3,\ 5)$ に $y=ax$ を\R{引}{ひ}きます。$S$ を $a$ の\R{式}{しき}で\R{表}{あらわ}して\R{平方完成}{へい|ほう|かん|せい}し、$S$ がいちばん\R{小}{ちい}さくなる $a$ を\R{求}{もと}めましょう。
\bkjans{$x^2$ の\R{和}{わ}14、$xy$ の\R{和}{わ}28、$y^2$ の\R{和}{わ}59。$S=14a^2-56a+59=14(a-2)^2+3$。$a=2$}
\end{jibun}
\end{minipage}\par
\endgroup

\begin{reidai}[\R{公式}{こう|しき}で $a$ を\R{求}{もと}めて、ばね\R{定数}{てい|すう}を\R{出}{だ}す]
ばねを\R{引}{ひ}く\R{力}{ちから} $x$（N）と、のび $y$（cm）をはかりました。\R{原点}{げん|てん}を\R{通}{とお}る\R{直線}{ちょく|せん} $y=ax$ を\R{最小二乗法}{さい|しょう|に|じょう|ほう}で\R{求}{もと}め、ばね\R{定数}{てい|すう} $k$（N/m）を\R{求}{もと}めましょう。
\begin{bkcenter}
\begin{bktable}{colspec={Q[l,m,wd=24mm] *{4}{Q[c,m,wd=10mm]}}, rowsep=0.5pt, cells={font=\small}}
\R{力}{ちから} $x$（N） & 1 & 2 & 3 & 4 \\
のび $y$（cm） & 2.2 & 3.9 & 6.0 & 8.0 \\
\end{bktable}
\end{bkcenter}
\kangaekata $a=$（$xy$ の\R{和}{わ}）÷（$x^2$ の\R{和}{わ}）で、$a$ は「1\,Nあたり\R{何}{なん}cmのびるか」です。ばね\R{定数}{てい|すう}は「1\,mのばすのに\R{何}{なん}Nいるか」なので、cmをmに\R{直}{なお}してから、わり\R{算}{ざん}の\R{向}{む}きを\R{逆}{ぎゃく}にします（2-15）。\par
\noindent\begin{minipage}[c]{0.62\linewidth}\raggedright
\shiki $x^2$ の\R{和}{わ}：$1+4+9+16=30$\par
\hspace*{2.6\zw}$xy$ の\R{和}{わ}：\par
\hspace*{2.6\zw}$2.2+7.8+18.0+32.0=60.0$\par
\hspace*{2.6\zw}$a=60.0\div30=2.0$\par
\hspace*{2.6\zw}1\,Nで\qty{2.0}{cm}のびる\par
\hspace*{2.6\zw}$\qty{2.0}{cm}\div\qty{100}{cm/m}=\qty{0.020}{m}$\par
\hspace*{2.6\zw}$k=\qty{1}{N}\div\qty{0.020}{m}=\qty{50}{N/m}$
\end{minipage}\hfill
\begin{minipage}[c]{0.37\linewidth}\centering
\begin{tikzpicture}
\begin{axis}[bkplot, width=46mm, height=44mm, xmin=0, xmax=4.6, ymin=0, ymax=9.5, xtick={0,...,4}, ytick={0,2,...,8}, xlabel={N}, ylabel={cm}]
  \addplot[bkBlue, line width=1.1pt, domain=0:4.4] {2*x};
  \addplot[only marks, mark=*, mark size=1.8pt, bkInk] coordinates {(1,2.2) (2,3.9) (3,6.0) (4,8.0)};
  \node[font=\sffamily\footnotesize, text=bkBlue!80!black, anchor=south east] at (axis cs:2.6,6.2) {$y=2.0x$};
\end{axis}
\end{tikzpicture}
\end{minipage}\par
\kotae $y=2.0x$、ばね\R{定数}{てい|すう}は\qty{50}{N/m}
\end{reidai}

\begin{jibun}
\R{力}{ちから} $x$（N）が1、2、3、4のとき、のび $y$（cm）は4.2、7.9、12.0、16.0でした。\R{原点}{げん|てん}を\R{通}{とお}る\R{直線}{ちょく|せん} $y=ax$ を\R{最小二乗法}{さい|しょう|に|じょう|ほう}で\R{求}{もと}め、ばね\R{定数}{てい|すう}を\R{求}{もと}めましょう。
\bkjans{$x^2$ の\R{和}{わ}30、$xy$ の\R{和}{わ}120.0。$a=120.0\div30=4.0$ で $y=4.0x$。$\qty{4.0}{cm}=\qty{0.040}{m}$、$k=1\div0.040=25$ で\qty{25}{N/m}}
\end{jibun}

\begin{reidai}[\R{直線}{ちょく|せん} $y=ax+b$ を、\R{共分散}{きょう|ぶん|さん}と\R{分散}{ぶん|さん}で\R{求}{もと}める]
あるばねの、おもりの\R{数}{かず} $x$ \R{個}{こ}と、\R{全体}{ぜん|たい}の\R{長}{なが}さ $y$〔cm〕です。\R{最小二乗法}{さい|しょう|に|じょう|ほう}で\R{直線}{ちょく|せん} $y=ax+b$ を\R{求}{もと}め、ばねのもとの\R{長}{なが}さを\R{答}{こた}えましょう。
\begin{bkcenter}
\begin{bktable}{colspec={Q[l,m,wd=24mm] *{5}{Q[c,m,wd=10mm]}}, rowsep=0.5pt, cells={font=\small}}
おもり $x$（\R{個}{こ}） & 1 & 2 & 3 & 4 & 5 \\
\R{長}{なが}さ $y$（cm） & 5.1 & 6.7 & 7.9 & 9.7 & 11.1 \\
\end{bktable}
\end{bkcenter}
\kangaekata \R{平均}{へい|きん}→\R{偏差}{へん|さ}→\R{偏差}{へん|さ}の\R{積}{せき}と $x$ の\R{偏差}{へん|さ}の2\R{乗}{じょう}→\R{共分散}{きょう|ぶん|さん}と\R{分散}{ぶん|さん}、の\R{順}{じゅん}に\R{計算}{けい|さん}します。$a=\sigma_{xy}\div\sigma_x^{\,2}$、$b=\bar{y}-a\bar{x}$ です。\par
\shiki $\bar{x}=15\div5=3$、$\bar{y}=40.5\div5=8.1$\par
\hspace*{2.6\zw}$x$ の\R{偏差}{へん|さ}：$-2$、$-1$、0、1、2\par
\hspace*{2.6\zw}$y$ の\R{偏差}{へん|さ}：$-3.0$、$-1.4$、$-0.2$、1.6、3.0\par
\hspace*{2.6\zw}\R{偏差}{へん|さ}の\R{積}{せき}：6.0、1.4、0、1.6、6.0で、\R{和}{わ}は15.0\par
\hspace*{2.6\zw}$x$ の\R{偏差}{へん|さ}の2\R{乗}{じょう}：4、1、0、1、4で、\R{和}{わ}は10\par
\hspace*{2.6\zw}$\sigma_{xy}=15.0\div5=3.0$、$\sigma_x^{\,2}=10\div5=2$、$a=3.0\div2=1.5$\par
\hspace*{2.6\zw}$b=8.1-1.5\times3=3.6$\par
\kotae $y=1.5x+3.6$。もとの\R{長}{なが}さ（$x=0$ のとき）は\R{約}{やく}\qty{3.6}{cm}
\end{reidai}

\begin{jibun}
おもり $x$ \R{個}{こ}が1〜5のとき、ばねの\R{長}{なが}さ $y$〔cm〕は14.0、15.9、18.1、20.1、21.9でした。\R{直線}{ちょく|せん} $y=ax+b$ を\R{最小二乗法}{さい|しょう|に|じょう|ほう}で\R{求}{もと}めましょう（③のばねC）。
\bkjans{$\bar{x}=3$、$\bar{y}=18.0$。\R{偏差}{へん|さ}の\R{積}{せき}の\R{和}{わ}20.0、$x$ の\R{偏差}{へん|さ}の2\R{乗}{じょう}の\R{和}{わ}10。$\sigma_{xy}=4$、$\sigma_x^{\,2}=2$、$a=2$、$b=18-6=12$。$y=2x+12$}
\end{jibun}

\begin{reidai}[\R{残差}{ざん|さ}を\R{計算}{けい|さん}して、\R{直線}{ちょく|せん}でよいか\R{調}{しら}べる]
5つの\R{点}{てん} $(1,\ 2.0)$、$(2,\ 2.5)$、$(3,\ 3.5)$、$(4,\ 5.0)$、$(5,\ 7.0)$ に、\R{最小二乗法}{さい|しょう|に|じょう|ほう}で\R{直線}{ちょく|せん}を\R{引}{ひ}くと $y=1.25x+0.25$ になりました。\R{残差}{ざん|さ}を\R{計算}{けい|さん}し、この\R{直線}{ちょく|せん}でよいかを\R{答}{こた}えましょう。\par
\kangaekata \R{残差}{ざん|さ}＝はかった\R{値}{あたい}－\R{線}{せん}の\R{値}{あたい}です。ならび\R{方}{かた}が「ばらばら」か「\R{山}{やま}や\R{谷}{たに}の\R{形}{かたち}」かを\R{見}{み}ます。\par
\shiki \R{線}{せん}の\R{値}{あたい}：$1.25\times1+0.25=1.5$、2.75、4.0、5.25、6.5\par
\hspace*{2.6\zw}\R{残差}{ざん|さ}：$2.0-1.5=0.5$、$2.5-2.75=-0.25$、$3.5-4.0=-0.5$、\par
\hspace*{2.6\zw}\phantom{\R{残差}{ざん|さ}：}$5.0-5.25=-0.25$、$7.0-6.5=0.5$\par
\kotae \R{残差}{ざん|さ}の\R{符号}{ふ|ごう}は＋、－、－、－、＋で、\R{谷}{たに}の\R{形}{かたち}。\R{直線}{ちょく|せん}は\R{合}{あ}っていない（\R{点}{てん}は\R{曲}{ま}がったきまりにしたがう）
\end{reidai}

\begin{reidai}[\R{直線}{ちょく|せん}で\R{予想}{よ|そう}する]
\bkreidairef{4}の\R{直線}{ちょく|せん} $y=1.5x+3.6$ で、（1）おもり3.5\R{個}{こ}\R{分}{ぶん}の\R{重}{おも}さをつるしたときの\R{長}{なが}さ、（2）おもり40\R{個}{こ}をつるしたときの\R{長}{なが}さを\R{予想}{よ|そう}しましょう。（2）の\R{予想}{よ|そう}は、\R{信用}{しん|よう}してよいですか。\par
\kangaekata $x$ に\R{代入}{だい|にゅう}します（2-13）。はかった\R{範囲}{はん|い}（1〜5\R{個}{こ}）の\R{外}{そと}の\R{予想}{よ|そう}は、\R{外挿}{がい|そう}といって\R{危}{あぶ}ないことを\R{思}{おも}い\R{出}{だ}しましょう（3-4）。\par
\shiki（1）\nobreak $y=1.5\times3.5+3.6=5.25+3.6=8.85$\par
\hspace*{2.6\zw}（2）\nobreak $y=1.5\times40+3.6=60+3.6=63.6$\par
\kotae（1）\nobreak \R{約}{やく}\qty{8.9}{cm}\quad（2）\nobreak \R{計算}{けい|さん}では\R{約}{やく}\qty{64}{cm}。でも、はかったのは5\R{個}{こ}までなので\R{信用}{しん|よう}できない。ばねは\R{引}{ひ}きすぎると、のびが\R{比例}{ひ|れい}しなくなる（3-0）
\end{reidai}

\begingroup\tcbset{equal height group=bkthrqjC}%
\noindent\begin{minipage}[t]{0.49\linewidth}
\begin{jibun}
5つの\R{点}{てん} $(1,\ 3.1)$、$(2,\ 4.8)$、$(3,\ 7.2)$、$(4,\ 8.8)$、$(5,\ 11.1)$ に、\R{最小二乗法}{さい|しょう|に|じょう|ほう}で\R{直線}{ちょく|せん} $y=2x+1$ を\R{引}{ひ}きました。\R{残差}{ざん|さ}を\R{計算}{けい|さん}し、この\R{直線}{ちょく|せん}でよいかを\R{答}{こた}えましょう。
\bkjans{\R{残差}{ざん|さ}は0.1、$-0.2$、0.2、$-0.2$、0.1。\R{小}{ちい}さく、＋と－がまざっている。\R{直線}{ちょく|せん}でよい}
\end{jibun}
\end{minipage}\hfill
\begin{minipage}[t]{0.49\linewidth}
\begin{jibun}
じぶんでとこう4の\R{直線}{ちょく|せん} $y=2x+12$ で、（1）おもり2.5\R{個}{こ}\R{分}{ぶん}、（2）おもり30\R{個}{こ}をつるしたときの\R{長}{なが}さを\R{予想}{よ|そう}しましょう。（2）は\R{信用}{しん|よう}してよいですか。
\bkjans{（1）\nobreak $2\times2.5+12=17$ で\R{約}{やく}\qty{17}{cm}\quad（2）\nobreak $2\times30+12=72$ で\R{約}{やく}\qty{72}{cm}。はかった\R{範囲}{はん|い}（1〜5\R{個}{こ}）の\R{外}{そと}なので\R{信用}{しん|よう}できない}
\end{jibun}
\end{minipage}\par
\endgroup

%% ---------------- まちがえやすいところ ----------------
\bksection{tsumazuki}{まちがえやすいところをチェック}

\begin{tsumazuki}{ずれの\R{和}{わ}が0なら、いちばんよい？}
\bkng{ずれをたすと0になったので、この\R{線}{せん}がいちばん\R{合}{あ}っている}
\bkok{ずれは2\R{乗}{じょう}してからたして、その\R{和}{わ}がいちばん\R{小}{ちい}さい\R{線}{せん}をさがす}
\noindent\begin{minipage}[c]{0.6\linewidth}
\R{上}{うえ}へのずれと\R{下}{した}へのずれは\R{打}{う}ち\R{消}{け}し\R{合}{あ}います。まったく\R{合}{あ}わない\R{線}{せん}でも、ずれの\R{和}{わ}が0になることがあります（\R{図}{ず}\ref{bk:3-16:fig-uchikeshi}）。\R{右}{みぎ}の\R{線}{せん}なら、ずれは $-12.5$ と $+12.5$ で\R{和}{わ}は0ですが、2\R{乗}{じょう}の\R{和}{わ}は $156.25+156.25=312.5$ と\R{大}{おお}きく、\R{合}{あ}っていないことがわかります。
\end{minipage}\hfill
\begin{minipage}[c]{0.36\linewidth}\centering
\begin{tikzpicture}[x=0.8mm, y=0.8mm]
  \draw[bkInk, line width=0.6pt, -{Stealth[length=3pt]}] (0,0) -- (56,0);
  \draw[bkInk, line width=0.6pt, -{Stealth[length=3pt]}] (0,0) -- (0,36);
  \draw[bkOrange, line width=1.2pt] (2,31) -- (54,5);
  \fill[bkInk] (16,11.5) circle (1.3) (40,24.5) circle (1.3);
  \draw[bkRed, line width=1.3pt] (16,11.5) -- (16,24);
  \draw[bkRed, line width=1.3pt] (40,24.5) -- (40,12);
  \node[font=\sffamily\footnotesize, text=bkRed!80!black, anchor=east] at (14.5,17.5) {$-12.5$};
  \node[font=\sffamily\footnotesize, text=bkRed!80!black, anchor=west] at (41.5,18) {$+12.5$};
\end{tikzpicture}
\end{minipage}
\end{tsumazuki}

\begin{tsumazuki}{$(y-ax)^2$ の\R{展開}{てん|かい}}
\bkng{$(y-ax)^2=y^2-a^2x^2$}
\bkok{$(y-ax)^2=y^2-2axy+a^2x^2$}
3-8の $(A-B)^2=A^2-2AB+B^2$ のまん\R{中}{なか}の $-2AB$ をわすれないこと。$y=4$、$a=3$、$x=1$ なら $(4-3)^2=1$ ですが、まちがいの\R{式}{しき}では $16-9=7$ になります。
\end{tsumazuki}

\begin{tsumazuki}{$a$ を「$y\div x$ の\R{平均}{へい|きん}」で\R{求}{もと}める}
\bkreidairef{2}の3つの\R{点}{てん}で\par
\bkng{$y\div x$ は4、3.5、2.67…なので、\R{平均}{へい|きん}して $a=3.39$}
\bkok{$a=$（$xy$ の\R{和}{わ}）÷（$x^2$ の\R{和}{わ}）$=42\div14=3$}
$a=3.39$ では $S=14\times0.39^2+3=5.1\cdots$ で、$a=3$ の $S=3$ より\R{大}{おお}きくなります。
\end{tsumazuki}

\begin{tsumazuki}{\R{傾}{かたむ}きを、\R{相関係数}{そう|かん|けい|すう}の\R{式}{しき}で\R{求}{もと}める}
\bkng{$a=\dfrac{\sigma_{xy}}{\sigma_x\,\sigma_y}$}
\bkok{$a=\dfrac{\sigma_{xy}}{\sigma_x^{\,2}}$（\R{分母}{ぶん|ぼ}は $x$ の\R{分散}{ぶん|さん}）}
\R{共分散}{きょう|ぶん|さん}を $\sigma_x\sigma_y$ でわると\R{相関係数}{そう|かん|けい|すう} $r$（3-15）で、\R{単位}{たん|い}のない\R{数}{かず}です。\R{傾}{かたむ}きには、cm/\R{個}{こ}のような\R{単位}{たん|い}があります。
\end{tsumazuki}


%% ---------------- 物理・工学 ----------------
\begin{butsuri}
\par\addvspace{3pt}\textbf{\R{計器}{けい|き}の\R{校正}{こう|せい}：\R{圧力}{あつ|りょく}をはかる\R{伝送器}{でん|そう|き}}\par
\noindent\begin{minipage}[c]{0.56\linewidth}
\R{工場}{こう|じょう}では、\R{配管}{はい|かん}の\R{中}{なか}の\R{圧力}{あつ|りょく}を「\R{伝送器}{でん|そう|き}」という\R{計器}{けい|き}ではかり、0〜\qty{100}{kPa}を4〜\qty{20}{mA}の\R{電流}{でん|りゅう}に\R{直}{なお}して、\R{遠}{とお}くの\R{制御室}{せい|ぎょ|しつ}へ\R{送}{おく}ります。\R{計器}{けい|き}が\R{正}{ただ}しいかを\R{確}{たし}かめるには、\R{正確}{せい|かく}な\R{圧力}{あつ|りょく}を\R{加}{くわ}えて、\R{出}{で}てくる\R{電流}{でん|りゅう}をはかります（\textbf{\R{校正}{こう|せい}}）。\bkidx{こうせい}{\R{校正}{こう|せい}}グラフの5\R{点}{てん}に\R{最小二乗法}{さい|しょう|に|じょう|ほう}を\R{使}{つか}うと、$y=0.15976x+4.014$ になりました。\R{決}{き}められた $y=0.16x+4$ とほぼ\R{同}{おな}じで、\R{残差}{ざん|さ}も\R{最大}{さい|だい}\qty{0.034}{mA}と\R{小}{ちい}さいので、\R{合格}{ごう|かく}です。\R{残差}{ざん|さ}が\R{山}{やま}や\R{谷}{たに}の\R{形}{かたち}なら、\R{計器}{けい|き}の\R{目}{め}もりが\R{直線}{ちょく|せん}からずれている（\R{直線性}{ちょく|せん|せい}が\R{悪}{わる}い）しるしです。
\end{minipage}\hfill
\begin{minipage}[c]{0.4\linewidth}\centering
\begin{tikzpicture}
\begin{axis}[bkplot, axis lines=left, width=50mm, height=46mm, xmin=0, xmax=110, ymin=0, ymax=22, xtick={0,50,100}, ytick={0,4,...,20}, xlabel={kPa}, ylabel={mA}]
  \addplot[bkBlue, line width=1pt, domain=0:105] {0.15976*x+4.014};
  \addplot[only marks, mark=*, mark size=1.7pt, bkRed] coordinates {(0,4.02) (25,8.01) (50,11.97) (75,16.03) (100,19.98)};
\end{axis}
\end{tikzpicture}
\end{minipage}

\par\addvspace{8pt}\textbf{\R{鉄}{てつ}のかたさ：ヤング\R{率}{りつ}}\par
\noindent \R{鉄}{てつ}の\R{棒}{ぼう}を\R{引}{ひ}っぱると、ほんの\R{少}{すこ}しのびます。「1\,mm$^2$あたりの\R{力}{ちから}（\R{応力}{おう|りょく}）」と、「もとの\R{長}{なが}さに\R{対}{たい}するのびの\R{割合}{わり|あい}（ひずみ）」は\R{比例}{ひ|れい}し、その\R{比例定数}{ひ|れい|てい|すう}を\textbf{ヤング\R{率}{りつ}}といいます。ひずみが0.0005、0.0010、0.0015、0.0020のとき、\R{応力}{おう|りょく}が101、199、302、\qty{398}{N/mm^2}だった\R{実験}{じっ|けん}に、\R{原点}{げん|てん}を\R{通}{とお}る\R{直線}{ちょく|せん}の\R{最小二乗法}{さい|しょう|に|じょう|ほう}を\R{使}{つか}うと、\R{傾}{かたむ}きは\R{約}{やく}\qty{200000}{N/mm^2}になります。\R{鉄}{てつ}のヤング\R{率}{りつ}として\R{知}{し}られている\R{値}{あたい}です。\R{橋}{はし}やビルの\R{設計}{せっ|けい}は、こうして\R{実験}{じっ|けん}から\R{決}{き}めた\R{数}{かず}をもとにしています。

\end{butsuri}
\begin{butsuri}
\par\addvspace{3pt}\textbf{\R{気温}{き|おん}から、あしたの\R{電気}{でん|き}を\R{予想}{よ|そう}する}\par
\noindent 3-15では、\R{夏}{なつ}の\R{気温}{き|おん}と\R{電気}{でん|き}の\R{使用量}{し|よう|りょう}に\R{正}{せい}の\R{相関}{そう|かん}があることを\R{見}{み}ました。ある\R{地域}{ち|いき}の10\R{日分}{にち|ぶん}のデータ（\R{最高気温}{さい|こう|き|おん}27〜36\,℃）に\R{最小二乗法}{さい|しょう|に|じょう|ほう}を\R{使}{つか}うと、「\R{最高気温}{さい|こう|き|おん}が1\,℃\R{上}{あ}がると、\R{最大電力}{さい|だい|でん|りょく}が\R{約}{やく}69\R{万}{まん}\,kWふえる」という\R{回帰直線}{かい|き|ちょく|せん}が\R{得}{え}られます。\R{天気予報}{てん|き|よ|ほう}の\R{気温}{き|おん}をこの\R{直線}{ちょく|せん}に\R{入}{い}れれば、あしたの\R{電気}{でん|き}の\R{使用量}{し|よう|りょう}を\R{予想}{よ|そう}して、\R{発電所}{はつ|でん|しょ}の\R{準備}{じゅん|び}ができるのです。
\end{butsuri}

%% ---------------- 章末問題 ----------------
\bkendproblems
\bkproblemgroup{\R{基本}{き|ほん}}{3}

\bkproblem 3つの\R{点}{てん} $(1,\ 1)$、$(2,\ 3)$、$(3,\ 4)$ があります。\R{直線}{ちょく|せん} $y=x$ と $y=1.5x$ のどちらがよく\R{合}{あ}っているかを、ずれの2\R{乗}{じょう}の\R{和}{わ}でくらべましょう。
\bkpans{$y=x$：ずれ0、1、1で $S=2$。$y=1.5x$：ずれ $-0.5$、0、$-0.5$ で $S=0.5$。$y=1.5x$ のほうがよく\R{合}{あ}う}

\bkproblem 3つの\R{点}{てん} $(1,\ 5)$、$(2,\ 9)$、$(3,\ 11)$ に $y=ax$ を\R{引}{ひ}きます。ずれの2\R{乗}{じょう}の\R{和}{わ} $S$ を $a$ の\R{式}{しき}で\R{表}{あらわ}して\R{平方完成}{へい|ほう|かん|せい}し、$S$ がいちばん\R{小}{ちい}さくなる $a$ を\R{求}{もと}めましょう。
\bkpans{$x^2$ の\R{和}{わ}14、$xy$ の\R{和}{わ}56、$y^2$ の\R{和}{わ}227。$S=14a^2-112a+227=14(a-4)^2+3$。$a=4$}

\bkproblem \R{力}{ちから} $x$（N）が1、2、3、4のとき、ばねののび $y$（cm）は1.2、1.9、3.0、4.0でした。\R{原点}{げん|てん}を\R{通}{とお}る\R{直線}{ちょく|せん} $y=ax$ を\R{最小二乗法}{さい|しょう|に|じょう|ほう}で\R{求}{もと}め、ばね\R{定数}{てい|すう}を\R{求}{もと}めましょう。
\bkpans{$x^2$ の\R{和}{わ}30、$xy$ の\R{和}{わ}30.0。$a=1.0$ で $y=1.0x$。$\qty{1.0}{cm}=\qty{0.010}{m}$、$k=1\div0.010=100$ で\qty{100}{N/m}}

\bkproblem 5つの\R{点}{てん} $(1,\ 2.1)$、$(2,\ 2.8)$、$(3,\ 4.2)$、$(4,\ 4.8)$、$(5,\ 6.1)$ に、\R{最小二乗法}{さい|しょう|に|じょう|ほう}で\R{直線}{ちょく|せん} $y=ax+b$ を\R{引}{ひ}きましょう。
\bkpans{$\bar{x}=3$、$\bar{y}=4.0$。\R{偏差}{へん|さ}の\R{積}{せき}の\R{和}{わ}10.0、$x$ の\R{偏差}{へん|さ}の2\R{乗}{じょう}の\R{和}{わ}10。$\sigma_{xy}=2$、$\sigma_x^{\,2}=2$、$a=1$、$b=4-3=1$。$y=x+1$}

\bkproblem 5つの\R{点}{てん} $(1,\ 0.5)$、$(2,\ 1.0)$、$(3,\ 2.0)$、$(4,\ 3.5)$、$(5,\ 5.5)$ に、\R{最小二乗法}{さい|しょう|に|じょう|ほう}で\R{直線}{ちょく|せん}を\R{引}{ひ}くと $y=1.25x-1.25$ になりました。\R{残差}{ざん|さ}を\R{計算}{けい|さん}し、この\R{直線}{ちょく|せん}でよいかを\R{答}{こた}えましょう。
\bkpans{\R{線}{せん}の\R{値}{あたい}は0、1.25、2.5、3.75、5.0。\R{残差}{ざん|さ}は0.5、$-0.25$、$-0.5$、$-0.25$、0.5で\R{谷}{たに}の\R{形}{かたち}。\R{直線}{ちょく|せん}は\R{合}{あ}っていない}

\bkproblem \R{問題}{もん|だい}4の\R{直線}{ちょく|せん}で、$x=2.5$ と $x=50$ のときの $y$ を\R{予想}{よ|そう}しましょう。$x=50$ の\R{予想}{よ|そう}は\R{信用}{しん|よう}してよいですか。
\bkpans{$x=2.5$ で3.5。$x=50$ で51。$x=50$ ははかった\R{範囲}{はん|い}（1〜5）の\R{外}{そと}なので、\R{信用}{しん|よう}できない}

\bkproblemgroup{チャレンジ}{3}

\bkproblem 3-0のソラさんのデータ（おもり0〜6\R{個}{こ}、のび0、1.6、2.9、4.6、5.9、7.6、8.9\,cm）で、\R{原点}{げん|てん}を\R{通}{とお}る\R{直線}{ちょく|せん} $y=ax$ を\R{最小二乗法}{さい|しょう|に|じょう|ほう}で\R{求}{もと}めましょう。おもり1\R{個}{こ}は\R{約}{やく}\qty{0.2}{N}として、ばね\R{定数}{てい|すう}を\R{求}{もと}めましょう。
\bkpans{$x^2$ の\R{和}{わ}91、$xy$ の\R{和}{わ}136.2。$a=136.2\div91=1.497$。$y=1.497x$。$\qty{1.497}{cm}=\qty{0.01497}{m}$、$k=0.2\div0.01497=13.4$ で\R{約}{やく}\qty{13}{N/m}}

\bkproblem \R{圧力計}{あつ|りょく|けい}の\R{校正}{こう|せい}で、0、50、\qty{100}{kPa}を\R{加}{くわ}えると、\R{出}{で}てくる\R{電流}{でん|りゅう}は4.0、12.1、\qty{19.9}{mA}でした。\R{最小二乗法}{さい|しょう|に|じょう|ほう}で\R{直線}{ちょく|せん} $y=ax+b$ を\R{求}{もと}めましょう（$x$ は kPa、$y$ は mA）。
\bkpans{$\bar{x}=50$、$\bar{y}=12.0$。\R{偏差}{へん|さ}の\R{積}{せき}の\R{和}{わ}は $(-50)\times(-8.0)+0\times0.1+50\times7.9=795$、$x$ の\R{偏差}{へん|さ}の2\R{乗}{じょう}の\R{和}{わ}5000。$a=795\div5000=0.159$、$b=12.0-0.159\times50=4.05$。$y=0.159x+4.05$}

\begin{furikaeri}
\begin{checklist}
\item ずれの2\R{乗}{じょう}の\R{和}{わ}で、どの\R{直線}{ちょく|せん}がよく\R{合}{あ}うかくらべられる。\bkfurref{bk:3-16:zure}
\item $y=ax$ の $a$ を、\R{平方完成}{へい|ほう|かん|せい}で\R{求}{もと}められる。\bkfurref{bk:3-16:genten}
\item $y=ax+b$ を\R{共分散}{きょう|ぶん|さん}と\R{分散}{ぶん|さん}で\R{求}{もと}め、\R{残差}{ざん|さ}で\R{確}{たし}かめられる。\bkfurref{bk:3-16:ippan}
\end{checklist}
\end{furikaeri}

\bkfukidashi{delta}{はかった\R{点}{てん}はいつもずれる。ずれの\R{正体}{しょう|たい}が「\R{偶然}{ぐう|ぜん}のばらつき」なら、どんな\R{形}{かたち}にちらばるのかな？\quad \R{次}{つぎ}の\R{章}{しょう}で\R{調}{しら}べよう。}
